This report presented some of the challenges involved in making a low-statistics ($\lesssim 1$ year) measurement of $\Delta m^2_{31}$ in JUNO. Statistical fluctuations can produce $\Delta\chi^2$ profiles that differ drastically from those obtained using Asimov data sets. Such fluctuations may lead to bad minima, meaning absolute minima that do not correspond to the true value used to generate the data set, as well as to multiple disjoint confidence intervals rather than the single, contiguous one normally expected for a measurement of this kind.

Moreover, the coverage obtained using the standard procedure based on Wilks' theorem, where confidence intervals are defined as regions within a constant value of $\Delta \chi^2$ (such as $\Delta \chi^2=1$ for the $1\sigma$ contour), was shown to be inadequate. This necessitates the implementation of an alternative approach, such as Feldman–Cousins, in which confidence intervals are derived from large ensembles of pseudo-experiments that include both statistical and systematic fluctuations. This, in turn, significantly increases the complexity and computational demands of the analysis.

It was demonstrated that these challenges affect almost exclusively the measurement of $\Delta m^2_{31}$ and not the solar parameters $\Delta m^2_{21}$ and $\sin^2\theta_{12}$, although care must still be taken to avoid biases associated with low-statistics bins. While not insurmountable, these challenges must be properly addressed to ensure an accurate determination of $\Delta m^2_{31}$. All this implies that the published $\Delta m^2_{31}$ sensitivities of Ref.~\cite{JUNO:2022osc}, which are based on Asimov data sets, have some limitations in that they do not capture all relevant statistical effects. In particular, low-statistics measurements of $\Delta m^2_{31}$ will not produce a single confidence interval matching the precision of the Asimov sensitivities with high probability. However, they remain largely accurate for larger exposures, corresponding to approximately one year of data or more.

This report also explored a potential mitigation strategy to avoid the occurrence of multiple disjoint intervals, which involves applying an external constraint on $\Delta m^2_{31}$. In this work, we examined the option of using a constraint from the Daya Bay experiment, although others could be used in principle. In this case, the external information helps the JUNO fit focus on the correct minimum, thereby drastically reducing the probability of obtaining multiple intervals. The drawback is that it makes the JUNO result dependent on external constraints.

The use of Daya Bay data to constrain the reactor antineutrino spectral shape was also explored, motivated by the fact that TAO data will likely not be available for the first low-statistics measurement of the oscillation parameters. This study presented a viable approach for incorporating the Daya Bay information through a joint fit with JUNO, starting from a common reactor antineutrino model that is allowed to vary. It was found that a precision comparable to that achievable with TAO can be obtained, with only a modest degradation in sensitivity observed for $\Delta m^2_{31}$. 

Finally, the potential impact of a low-statistics JUNO measurement of $\Delta m^2_{31}$ on the mass ordering determination of ongoing long-baseline accelerator experiments was also investigated. As expected, the extent of this impact depends primarily on the best-fit value and the precision achieved by JUNO for $\Delta m^2_{31}$. In some scenarios, the improvement can be substantial, potentially enabling a $3\sigma$ determination of the mass ordering.